% SGA 3, Expose VI_A, Loop 1 English reconstruction.
% Scope: paragraph 4.2 in full; paragraph 4.3 is excluded.
% Authority: Polo--Gille Exp6A-13oct24.pdf, component-local page 27
% (combined-reader page 330; printed page 319).
% OCR and the Reinhold English Markdown are comparison witnesses only.

\medskip
\phantomsection\label{par:sga3-VIA-4.2}%
\noindent\textbf{4.2.---}
特别地, 取 \(G^\alpha\) 为 \(G\) 中单位元所在的连通分支 \(G^0\).
于是 \(\mathscr S(G^0)\) 显然包含 \(F\) 在 \(u\) 下的像,
后者恰为单位元的等价类.
另一方面, 若 \(F^\beta\) 为 \(F\) 的连通分支,
则由\hyperref[lem:sga3-VIA-2.1.2]{2.1.2}, \(F^\beta\times_A G^0\) 连通,
故其在 \(\lambda\) 下的像包含于 \(G\) 中 \(u(F^\beta)\) 所在的连通分支.
换言之, \(\mathscr S(G^0)\) 为与 \(F\) 的像相交的连通分支之并.

还须注意, 底空间为 \(\mathscr S(G^0)\) 的 \(G\) 的开子概形是 \(G\) 的子群;
仍用 \(\mathscr S(G^0)\) 表示它.
事实上, \(G\) 上的取逆映射保持 \(F\) 的像,
且置换 \(G\) 中与此像相交的连通分支. 因此只需证明
\[
  \nu:G\times_A G\longrightarrow G
\]
将 \(\mathscr S(G^0)\times_A\mathscr S(G^0)\) 映入 \(\mathscr S(G^0)\).
为此可假设 \(A\) 为代数闭域:
采用\hyperref[par:sga3-VIA-4.1]{4.1} 的记号,
\(\mathscr S(G^0)\otimes_A\bar k\) 是 \((G\otimes_A\bar k)^0\)
对同态 \(u\otimes_A\bar k\) 所定义的等价关系的饱和.

若 \(G^\gamma\) 和 \(G^\delta\) 为 \(\mathscr S(G^0)\) 的连通分支,
则 \(G^\gamma\times_A G^\delta\) 连通, 且其在 \(\nu\) 下的像与 \(F\) 的像相交.
因此,
\[
  \nu(G^\gamma\times_A G^\delta)
\]
包含于 \(G\) 中一个与 \(u(F)\) 相交的连通分支.\begingroup
\renewcommand{\thefootnote}{\fnsymbol{footnote}}%
\footnote[1]{英译者注: 法文底本在此末句中印作 \(u(G^\gamma\times_A G^\delta)\).
本段定义并始终使用的映射是乘法 \(\nu:G\times_A G\to G\);
因此译文采用 \(\nu\).}%
\endgroup
% Independent prose review: source and translation compared paragraph by paragraph through this component; statements, hypotheses, proofs, notes and diagram captions covered; bibliography reserved for its assigned reviewer; source math and reference keys unchanged.
